\chapter{Conclusiones}
% BORRADOR autónomo (2026-08-08). Editar a tu voz. UAX: empezar por objetivos específicos y
% terminar por el general, mostrando su cumplimiento. Recomendado: 5-10 páginas.

Este capítulo recupera los objetivos planteados y valora su cumplimiento, recorre las limitaciones
y la prospectiva, y cierra con una reflexión personal.

Respecto a los \textbf{objetivos específicos}, todos se han cumplido. Se ha \textbf{revisado el
estado del arte} (Cap.~3), situando la aportación en el hueco de la explicabilidad validada de los
MLLM de radiografía de tórax. Se ha \textbf{diseñado e implementado} un generador multimodal propio
con destilación de conocimiento y un marco reutilizable de explicabilidad. Se ha \textbf{cuantificado
de forma objetiva el anclaje clínico} de la atención contra anotaciones reales de localización
(lift global 3,03). Y se ha \textbf{determinado la capacidad de la explicabilidad y la incertidumbre
para detectar errores}, con un resultado matizado y sólido: la morfología de la explicación no
predice la fiabilidad —lo que se confirma por triple vía, incluidas las cajas de anclaje del propio
modelo—, pero la \textbf{incertidumbre de la decisión detecta las alucinaciones} (AUC 0,90 tras
depurar el token de apertura), no las omisiones, que solo se captan de forma parcial.

En cuanto al \textbf{objetivo general} —evaluar la fidelidad y el anclaje de las explicaciones y en
qué medida permiten estimar la fiabilidad—, se ha alcanzado y ampliado: las explicaciones son fieles
y ancladas, pero la fiabilidad reside en la incertidumbre de la decisión, no en el mapa de calor. Es
más, dicha señal ha demostrado ser \textbf{agnóstica al modelo y robusta al dominio}: validada en un
diseño \(2\times2\) (dos generadores \(\times\) dos dominios), se transfiere a un segundo modelo del
estado del arte (MAIRA-2) y se mantiene \emph{fuera de distribución} para el generador de calidad
suficiente (AUC 0,80--0,81), mientras que la degradación observada se atribuye al generador débil, no
al método. La contribución es, por tanto, triple: la \textbf{auditoría objetiva del anclaje} de un
MLLM de CXR; la \textbf{disociación} entre errores detectables (fabricaciones) e indetectables por
confianza (omisiones); y una \textbf{señal de confianza simple, agnóstica, calibrada y validada de
forma cruzada}, junto con su traducción clínica en una demostración interactiva. Como aportación
metodológica menor pero de valor práctico, se documenta que \emph{excluir el token de apertura del
informe} mejora de forma sistemática la señal de incertidumbre.

El significado de conjunto de estas conclusiones puede resumirse en una idea: en los modelos
multimodales de radiografía de tórax, la fiabilidad de un informe no se lee en \emph{lo que el modelo
muestra} —su explicación visual— sino en \emph{lo seguro que está} al decidir, y esa seguridad es
cuantificable, transferible entre modelos y calibrable para su uso clínico, siempre que se comunique
con honestidad su alcance. Este desplazamiento del foco —de la explicación a la incertidumbre— es la
contribución conceptual central del trabajo, y tiene implicaciones directas para el diseño de sistemas
de ayuda a la decisión: la evidencia visual debe presentarse como \emph{apoyo a la localización}, no
como garantía de acierto, y la confianza debe construirse sobre la señal interna del modelo, no sobre
la forma de su explicación.

\clearpage
\section{Respuesta a las preguntas de investigación}

\begin{table}[h]\centering
\caption{Resumen de las cuatro preguntas de investigación (Capítulo~1) e hipótesis asociadas
(H1--H4, Sección~4.2) frente a la evidencia obtenida.}
\label{tab:resumen_pi}
\begin{tabular}{lp{2.6cm}p{5.5cm}}
\toprule
Pregunta & Resultado & Evidencia clave \\ \midrule
PI1 — ¿anclaje clínico? & Confirmada & Lift global 3,03; \textit{pointing-game} 0,38 \\
PI2 — ¿morfología predice fiabilidad? & Refutada & AUC 0,54 (n.s.); confirmado por triple vía \\
PI3 — ¿incertidumbre detecta errores? & Confirmada (asimétrica) & AUC 0,80--0,90 alucinaciones; parcial en omisiones \\
PI4 — ¿agnóstica y robusta? & Confirmada con matiz & Transfiere a MAIRA-2; falla solo en generador débil OOD \\
\bottomrule
\end{tabular}
\end{table}

Recapitulando las cuatro preguntas planteadas en el Capítulo~1: en cuanto a \textbf{PI1}, la atención
está clínicamente anclada (lift 3,03; tres veces el azar en el \textit{pointing-game}), lo que valida
su uso como puntero de \emph{dónde} mira el modelo. Respecto a \textbf{PI2}, la morfología de la
explicación \emph{no} predice la fiabilidad: ni el mapa de calor de atención, ni el \textit{lift}
objetivo, ni las cajas de anclaje entrenadas del propio modelo —que localizan las alucinaciones con
la misma plausibilidad que los hallazgos reales— separan el acierto del error. En cuanto a
\textbf{PI3}, la incertidumbre de la decisión sí detecta los errores, pero de forma \emph{asimétrica}:
robusta para las alucinaciones (AUC 0,80--0,90) y solo parcial para las omisiones, que tienden a
situarse cerca del umbral. Y sobre \textbf{PI4}, la señal es agnóstica al modelo y robusta al dominio:
se transfiere a un segundo generador y se mantiene fuera de distribución para un generador de calidad
suficiente, degradándose solo con un generador débil que, sin datos del dominio, alucina con
confianza. En conjunto, las cuatro preguntas quedan respondidas de forma coherente y con matices
respaldados por la evidencia.

\section{Limitaciones}
Algunas limitaciones de las versiones preliminares se han \emph{resuelto} en el propio trabajo: la
dependencia de un único juez se ha mitigado con \textbf{tres fuentes de verdad} concordantes; la de un
único modelo, validando en \textbf{dos generadores}; y la de un único dominio, con la \textbf{réplica
en PadChest-GR}. Persisten, no obstante, limitaciones reales. El número de alucinaciones por celda es
\textbf{reducido} (entre 12 y 110), lo que ensancha los intervalos de confianza. La evaluación
\emph{in-distribution} de MAIRA-2 sobre PadChest puede incluir \textbf{solape train--test}, al haber
sido entrenado con ese corpus; se mitiga usando Open-I —ajeno a su entrenamiento— como prueba fuera de
distribución, pero conviene señalarlo. La detección de \textbf{omisiones} es solo parcial: son errores
de sensibilidad del generador, no siempre acompañados de incertidumbre. Metodológicamente, varias
señales de forma resultaron engañosas (el máximo de atención, la entropía confundida con la longitud,
la relación \emph{lift}--acierto), lo que obligó a un control cuidadoso de confusores. Por último, el
\textbf{beneficio clínico no se ha demostrado empíricamente}: requeriría un estudio prospectivo de
lectura.

Varias de estas limitaciones tienen origen en las condiciones materiales del proyecto. El
\textbf{acceso a datos} clínicos de mayor tamaño estuvo restringido: MIMIC-CXR, el corpus narrativo
de referencia en el área, exige credenciales de PhysioNet no disponibles para este trabajo
(Sección~4.1), lo que obligó a construir el diseño experimental sobre conjuntos abiertos de menor
tamaño (Open-I, NIH ChestX-ray8, PadChest-GR) y es, en buena medida, la causa directa del reducido
número de alucinaciones por celda señalado más arriba. El \textbf{cómputo} también estuvo acotado a
una única GPU de consumo, lo que impuso ejecutar los modelos del estado del arte cuantizados a 8
bits —sin comparación frente a precisión completa— y limitó la escala y el número de repeticiones
de las ablaciones del modelo propio (Sección~4.4). Ninguna de estas restricciones invalida los
resultados, que se sostienen sobre tres fuentes de verdad y dos réplicas independientes, pero sí
marca el techo de precisión alcanzable con los recursos disponibles y es la razón principal de la
primera línea de prospectiva (consolidación estadística).

\section{Prospectiva}
El trabajo abre varias líneas de continuación, que se ordenan de la más inmediata a la más ambiciosa.

\textbf{Consolidación estadística.} La primera línea es \textbf{ampliar la muestra} de imágenes por
celda del diseño, dado que el número de alucinaciones —el factor que limita la precisión de la
estimación— es reducido; ampliarlo estrecharía los intervalos de confianza. En la misma dirección,
añadir más generadores (por ejemplo, modelos con razonamiento por agentes o de otra familia) y más
dominios reforzaría la generalidad de la conclusión sobre la robustez de la señal.

\textbf{Cierre de la duda del solape.} La evaluación \emph{in-distribution} de MAIRA-2 sobre PadChest
podría incluir solape con su entrenamiento; evaluar el modelo fuera de distribución en dominios
verdaderamente ajenos y sin intersección con su corpus cerraría definitivamente esta cuestión.

\textbf{Detección de omisiones.} Las omisiones son el error que la incertidumbre no capta de forma
fiable, por tratarse de fallos confiados de sensibilidad. Abordarlas requeriría vías complementarias:
mejorar la sensibilidad del propio generador, emplear un anclaje explícito que ``obligue'' a mirar
regiones sospechosas, o recurrir a la \textbf{auto-consistencia semántica} entre múltiples muestras,
que podría revelar la inestabilidad de un informe incompleto.

\textbf{De la detección a la prevención.} Más allá de detectar la alucinación tras generarla, la señal
de incertidumbre podría integrarse en el propio proceso de \emph{decodificación} para \textbf{prevenir}
la fabricación —por ejemplo, penalizando las continuaciones de baja confianza o solicitando evidencia
adicional—, convirtiendo el sistema de un detector \emph{a posteriori} en una salvaguarda \emph{a
priori}.

\textbf{Validación clínica.} Finalmente, la línea de mayor calado es un \textbf{estudio de lectura}
prospectivo con radiólogos que mida el beneficio real de presentar informe, anclaje y confianza,
cuantificando su efecto sobre la precisión diagnóstica, el tiempo de lectura y el riesgo de sesgo de
automatización. Solo un estudio de este tipo puede demostrar el valor clínico que este trabajo
fundamenta pero no llega a probar, y sería el paso natural hacia una eventual integración del sistema
como módulo asistivo en el flujo de trabajo radiológico.

\clearpage
\section{Mejoras futuras del sistema}
Las líneas anteriores extienden la investigación; las que siguen extienden el \textbf{sistema} en un
sentido de ingeniería, orientado a acercarlo a un entorno de uso real.

\textbf{Ampliación del catálogo de hallazgos.} El etiquetado objetivo se restringe a los hallazgos
\emph{agudos} más frecuentes y a las ocho patologías con caja de NIH ChestX-ray8. Extender el marco a
un catálogo más amplio —hallazgos crónicos, óseos o de dispositivos— exigiría nuevas fuentes de verdad,
pero ampliaría la cobertura clínica real del sistema.

\textbf{Integración con el flujo de trabajo radiológico.} La demostración actual es un prototipo
\textit{standalone}; su valor práctico aumentaría integrándola como módulo dentro del visor de
imágenes o del sistema de información radiológica (RIS/PACS) del centro, de modo que el informe, el
anclaje y la confianza aparezcan directamente en la estación de lectura, sin cambiar la herramienta que
el radiólogo ya usa.

\textbf{Optimización de la latencia de inferencia.} El sistema se ejecuta hoy cuantizado a 8 bits en
una única GPU de consumo (Sección~4.7); técnicas adicionales de optimización —cuantización a 4 bits,
motores de inferencia especializados o destilación del propio generador— reducirían la latencia y el
coste de despliegue, acercando el sistema a un uso en tiempo real durante la lectura.

\textbf{Ampliación a otras modalidades.} El marco de explicabilidad e incertidumbre —agnóstico al
generador por construcción— no depende de particularidades de la radiografía de tórax; extenderlo a
otras modalidades con MLLM propios (tomografía computarizada, resonancia magnética, mamografía)
permitiría comprobar si la conclusión central del trabajo —que la fiabilidad reside en la incertidumbre
y no en la explicación visual— se sostiene fuera del dominio aquí estudiado.

\textbf{Modelos de mayor escala y capacidad de imagen.} Este trabajo se ha ejecutado sobre generadores
de 7 mil millones de parámetros y con la resolución de entrada que cada uno define por diseño,
condicionado por el hardware de consumo disponible (Sección~6.2). Investigar con modelos de mayor
potencia y con mayor capacidad de procesamiento de imagen —más tokens visuales, mayor resolución de
entrada— permitiría estudiar si la morfología del mapa de calor, que en este trabajo no logró predecir
la fiabilidad, se comporta de forma distinta en arquitecturas más grandes y con una representación
visual más rica, cerrando una pregunta que la limitación de recursos deja abierta.

\textbf{Consolidación de la portabilidad.} El trabajo ya distribuye la memoria y las demostraciones
mediante imágenes \textbf{Docker} reproducibles (Anexo de organización del repositorio); una línea
natural es exponer el sistema de análisis como una \textbf{API} desplegable en la nube, desacoplada de
la interfaz Gradio actual, para que otros equipos puedan integrarlo en sus propias herramientas sin
depender de esta implementación concreta.

\section{Consideraciones finales}
Considero que el resultado final de esta investigación fue satisfactorio, no solo por lo logrado al
poder ``acotar'' con una métrica un umbral de alucinaciones —con el beneficio añadido de que muchas
omisiones se sitúan precisamente en ese umbral—, sino por todo el aprendizaje que ha conllevado.

Ese aprendizaje me ha enseñado, incluso creyendo que ya estaba familiarizado con el método científico,
lo que de verdad significa el rigor en el control de confusores tras varios ``falsos positivos''.
Cuántas veces habré creído tener algo interesante entre manos, para que la siguiente batería de
pruebas mostrara que la hipótesis no se sostenía. Pero esa negación de la idea inicial del TFM no fue
el final del proyecto, sino el paso de una intuición —los mapas de calor— a su refutación con
evidencia. He aprendido a trabajar con MLLM de gran escala, con cuantización y con extracción de
señales a escala; a comunicar un resultado negativo con la misma solidez que uno positivo; y,
finalmente, a llegar a algo que sí era capaz de aportar.

Si hoy empezara de nuevo —no con el conocimiento adquirido, sino con las aptitudes desarrolladas
durante el Máster y el propio TFM— seguramente no habría cambiado el resultado, pues se habrían hecho
las mismas pruebas, pero las habría hecho bien a la primera, con más rigor científico y sin dejarme
llevar por pequeñas victorias sesgadas que acabaron en sobretrabajo antes de comprobar si realmente
eran buenos indicios. Y, sobre todo, habría empezado el procesamiento de datos mucho antes, incluso
mientras todavía planteaba las ideas. Me habría encantado poder validar las hipótesis con más datos,
pero el tiempo disponible y el hardware del equipo con el que se realizó el trabajo lo complicaron.

No puedo terminar este trabajo sin expresar mi deseo de que la asistencia por IA en radiología
evolucione hacia sistemas honestos sobre su propia incertidumbre, y no solo hacia modelos cada vez más
grandes —y, en muchos casos, aún más caja negra—, y de que se apliquen a los MLLM todos los avances de
la IA explicable clásica, como los vistos en la asignatura de XAI del Máster. Espero que esta pequeña
contribución sea un paso, por modesto que sea, en esa dirección.
